\scanpage{52}
% archive sheet number: 51
\margin{\uline{! : \unsure{cas.} affine} \deleted{au \unsure{sens} \unsure{syst. proj.} de \uuline{\unsure{groupes}}}}
\diagram{\begin{tikzcd}[ampersand replacement=\&,column sep=large,row sep=large]
\operatorname{Progr}^{!}(C'_0)
  \arrow[r,"\deleted{\unsure{\leftarrow}}","\alpha"']
\& \operatorname{Gr}(C)
  \arrow[r,phantom,"{\underset{\beta}{\overset{\text{\unsure{équiv.}}}{\simeq}}}" description]
\& \operatorname{Gr}(C')
  \arrow[r,"{\substack{\text{pl. fid}\\\deleted{\unsure{\operatorname{Gr}(C)}}}}","\unsure{\sigma|\tau}"']
\& \operatorname{Gr}\ \text{qucpt}\ (\mathrm{Sch}_{/k})
  \arrow[r,"{\substack{\deleted{\text{pl. fid}}\\\unsure{\ell}}}"]
\& {} \\
\operatorname{Progr}(C'_0)\arrow[u,hook,"i\ \text{pl. fid}"]
  \arrow[rr,"\unsure{\kappa|k}"']
\& \& \operatorname{Gr}\operatorname{Pro}(C'_0)
  \arrow[u,hook,"j\ \text{pl. fid}"']
\end{tikzcd}}

$\beta$ \unsure{équivalence} car $\deleted{\unsure{\operatorname{Gr}(C)\simeq\operatorname{Gr}(}}\ C\simeq C'$.

$\sigma$ pl. fid \unsure{trivial} sur $\operatorname{Gr}(C)\longrightarrow\operatorname{Gr}\ \text{qu. cpts}\ (\mathrm{Sch}_{/k})$.

$i,j$ pl. fid \deleted{\unsure{est}} \unsure{trivial} [N.B $C'\unsure{-|\longrightarrow}\operatorname{Pro}(C'_0)$ est pl. fid et \unsure{commute} aux $\varprojlim$ \unsure{finies}]

\[
\tikz[remember picture,baseline=(fid52.base)]{\node[inner sep=2pt] (fid52) {$\kappa\ \text{pl. fid}$};\draw (fid52.center) ellipse (0.85cm and 0.35cm);}\ \Longrightarrow\ \alpha\ \text{pl. fid}\ \Longleftrightarrow\
[\sigma\alpha:\operatorname{Progr}^{!}(C'_0)\longrightarrow\operatorname{Gr}\ \text{qucpt}\ (\mathrm{Sch}_{/k})\ \deleted{\text{\unsure{est}}}\ \text{pl. fid.}]
\]
\margin{\tikz[remember picture,baseline=(vrai52.base)]{\node[draw,inner sep=2pt] (vrai52) {vrai};\draw[overlay] (vrai52.north) -- (fid52.south);}}
% The boxed "vrai" is joined by a line to the circled assertion kappa pl. fid.
\diagram{\begin{tikzcd}[ampersand replacement=\&,column sep=large,row sep=large,remember picture]
\sigma\alpha\ \text{une \unsure{équivalence}}
  \arrow[r,Leftrightarrow]
\& \sigma\ \text{et}\ \alpha\ \text{des \unsure{équiv.}}
  \arrow[d,Rightarrow]\arrow[r,Rightarrow]
\& \begin{gathered}\deleted{\text{\unsure{tout}}}\ \text{schéma en}\\
\text{\unsure{gr.} qu. cpt sur}\\ k\ \text{admet un}\\
\text{\unsure{noyau}}\ \text{\uline{affine}}\end{gathered} \\
|[alias=eq52]|\kappa\ \text{une \unsure{équivalence}}\arrow[r,Rightarrow,""{name=cond52}]
\& \alpha\ \text{une \unsure{équivalence}}
\& \begin{gathered}\text{donc un schéma en}\\
\text{\unsure{gr. quot} de t. f. sur }k\end{gathered}
  \arrow[ul,Rightarrow,bend left=25,"\text{\unsure{un peu de travail}}"']
\end{tikzcd}}
\margin{$\unsure{\kappa|K}$}
\margin{\tikz[remember picture,baseline=(d52.base)]{\node[inner sep=2pt] (d52) {\unsure{douteux}};\draw (d52.center) ellipse (0.85cm and 0.3cm);\draw[overlay] (d52.north) -- (eq52.south);}}
% The circled "douteux" is linked to "kappa une équivalence".
\margin{\tikz[remember picture,baseline=(note52.base)]{\node[inner sep=0pt,text width=2.2cm,align=left] (note52) {\unsure{pourvu qu'on} \deleted{\unsure{sache}} \unsure{qu'un|puisse} \unsure{un} \unsure{pro. gr.} tel que le \unsure{proschéma} \unsure{associé} \unsure{soit} \unsure{ess.|un.} affine, \unsure{soit} \unsure{lui-même} \unsure{ess.|un sch.} affine.};\draw[overlay] (note52.north) to[out=90,in=-90] (cond52.south);}}
% This last note is connected to the implication kappa equivalence => alpha equivalence.
